% This LaTeX file was created by Metamath on 22-Jul-2024 3:29 PM.
\begin{restatable}{metadefinition}{dfblat}~\blTitle{df-bl.at}~ Definition of @ for predicates.
   @ determines in which world to evaluate.
\begin{align}
\vdash \tau \in \mathbf{T}\quad\Rightarrow\quad \vdash ( \mathbf{I} \lsem (
    \mw{\varphi} @ \tau ) \rsem \leftrightarrow \mathbf{I} \lsem \mw{\varphi} ,
    \tau \rsem )\label{eq:df-bl.at}\tag{df-bl.at}
\end{align}
\end{restatable}

\begin{restatable}{metadefinition}{dfblatc}~\blTitle{df-bl.atc}~ Definition of @ for values.
   @ determines in which world to evaluate.
\begin{align}
\vdash \tau \in \mathbf{T}\quad\Rightarrow\quad \vdash \mathbf{I} \lsem (
    \mc{A} @ \tau ) \rsem = \mathbf{I} \lsem \mc{A} , \tau
    \rsem\label{eq:df-bl.atc}\tag{df-bl.atc}
\end{align}
\end{restatable}

\begin{restatable}{metadefinition}{dfblat2}~\blTitle{df-bl.at2}~ Definition of Double @ for predicates.
   Second application of @ is inconsequential, predicate.
\begin{align}
\vdash \tau_1 \in \mathbf{T}\quad\&\quad \vdash \tau_2 \in
    \mathbf{T}\quad\Rightarrow\quad \vdash ( \mathbf{I} \lsem ( ( \mw{\varphi}
    @ \tau_1 ) @ \tau_2 ) \rsem \leftrightarrow \mathbf{I} \lsem \mw{\varphi} ,
    \tau_1 \rsem )\label{eq:df-bl.at2}\tag{df-bl.at2}
\end{align}
\end{restatable}

\begin{restatable}{metadefinition}{dfblat2c}~\blTitle{df-bl.at2c}~ Definition of Double @ for values.
   Second application of @ is inconsequential, value.
\begin{align}
\vdash \tau_1 \in \mathbf{T}\quad\&\quad \vdash \tau_2 \in
    \mathbf{T}\quad\Rightarrow\quad \vdash \mathbf{I} \lsem ( ( \mc{A} @ \tau_1
    ) @ \tau_2 ) \rsem = \mathbf{I} \lsem \mc{A} , \tau_1
    \rsem\label{eq:df-bl.at2c}\tag{df-bl.at2c}
\end{align}
\end{restatable}

\begin{restatable}{lemma}{blatintro}~\blTitle{bl.atintro}~ @ Introduction for predicates.
   If a predicate is always true it's true at any time.
\begin{align}
\vdash \tau \in \mathbf{T}\quad\Rightarrow\quad \vdash ( \mathbf{I} \lsem
    \mw{\varphi} \rsem \rightarrow \mathbf{I} \lsem ( \mw{\varphi} @ \tau )
    \rsem )\label{eq:bl.atintro}\tag{bl.atintro}
\end{align}
\end{restatable}

\begin{proof}
\begin{align}
  1 &&  & \vdash \tau \in \mathbf{T}&\text{Hyp~atintro.1}\notag%bl.atintro.1
  \\ 2 &&  & \vdash ( \mathbf{I} \lsem \mw{\varphi} \rsem \rightarrow
      \mathbf{I} \lsem \mw{\varphi} , \tau \rsem
      )&(\mbox{\ref{eq:ax-bl.models}},1)\notag
  \\ 3 &&  & \vdash \tau \in \mathbf{T}&\text{Hyp~atintro.1}\notag%bl.atintro.1
  \\ 4 &&  & \vdash ( \mathbf{I} \lsem ( \mw{\varphi} @ \tau ) \rsem
      \leftrightarrow \mathbf{I} \lsem \mw{\varphi} , \tau
      \rsem )&(\mbox{\ref{eq:df-bl.at}},3)\notag
  \\ 5 &&  & \vdash ( \mathbf{I} \lsem \mw{\varphi} \rsem \rightarrow
      \mathbf{I} \lsem ( \mw{\varphi} @ \tau ) \rsem
      )&(\mbox{\ref{eq:sylibr}},2,4)\notag
\end{align}
\end{proof}

\begin{restatable}{metadefinition}{dfblatbi}~\blTitle{df-bl.atbi}~ Definition of @ Equivalence
(biconditional) for predicates.
\begin{align}
\vdash \tau_1 \in \mathbf{T}\quad\&\quad \vdash \tau_2 \in
    \mathbf{T}\quad\Rightarrow\quad \vdash ( ( \tau_1 = \tau_2 \wedge (
    \mw{\varphi} \leftrightarrow \mw{\psi} ) ) \rightarrow ( \mathbf{I} \lsem (
    \mw{\varphi} @ \tau_1 ) \rsem \leftrightarrow \mathbf{I} \lsem ( \mw{\psi}
    @ \tau_2 ) \rsem ) )\label{eq:df-bl.atbi}\tag{df-bl.atbi}
\end{align}
\end{restatable}

\begin{restatable}{lemma}{blatbii}~\blTitle{bl.atbii}~ @ Equivalence (biconditional) for
predicates, inference.
\begin{align}
\vdash \tau_1 \in \mathbf{T}\quad\&\quad \vdash \tau_2 \in
    \mathbf{T}\quad\&\quad \vdash ( \mw{\varphi} \leftrightarrow \mw{\psi}
    )\quad\&\quad \vdash \tau_1 = \tau_2\quad\Rightarrow\quad \vdash (
    \mathbf{I} \lsem ( \mw{\varphi} @ \tau_1 ) \rsem \leftrightarrow \mathbf{I}
    \lsem ( \mw{\psi} @ \tau_2 ) \rsem )\label{eq:bl.atbii}\tag{bl.atbii}
\end{align}
\end{restatable}

\begin{proof}
\begin{align}
  1 &&  & \vdash \tau_1 = \tau_2&\text{Hyp~atbii.4}\notag%bl.atbii.4
  \\ 2 &&  & \vdash ( \mw{\varphi} \leftrightarrow \mw{\psi}
      )&\text{Hyp~atbii.3}\notag%bl.atbii.3
  \\ 3 &&  & \vdash \tau_1 \in \mathbf{T}&\text{Hyp~atbii.1}\notag%bl.atbii.1
  \\ 4 &&  & \vdash \tau_2 \in \mathbf{T}&\text{Hyp~atbii.2}\notag%bl.atbii.2
  \\ 5 &&  & \vdash ( ( \tau_1 = \tau_2 \wedge ( \mw{\varphi} \leftrightarrow
      \mw{\psi} ) ) \rightarrow ( \mathbf{I} \lsem (
      \mw{\varphi} @ \tau_1 ) \rsem \leftrightarrow
      \mathbf{I} \lsem ( \mw{\psi} @ \tau_2 ) \rsem )
      )&(\mbox{\ref{eq:df-bl.atbi}},3,4)\notag
  \\ 6 &&  & \vdash ( \mathbf{I} \lsem ( \mw{\varphi} @ \tau_1 ) \rsem
      \leftrightarrow \mathbf{I} \lsem ( \mw{\psi} @ \tau_2
      ) \rsem )&(\mbox{\ref{eq:mp2an}},1,2,5)\notag
\end{align}
\end{proof}

\begin{restatable}{metadefinition}{dfblateqc}~\blTitle{df-bl.ateqc}~ Definition of @ Equality for values.
\begin{align}
\vdash \tau_1 \in \mathbf{T}\quad\&\quad \vdash \tau_2 \in
    \mathbf{T}\quad\Rightarrow\quad \vdash ( ( \tau_1 = \tau_2 \wedge \mc{A} =
    \mc{B} ) \rightarrow \mathbf{I} \lsem ( \mc{A} @ \tau_1 ) \rsem =
    \mathbf{I} \lsem ( \mc{B} @ \tau_2 ) \rsem
    )\label{eq:df-bl.ateqc}\tag{df-bl.ateqc}
\end{align}
\end{restatable}

\begin{restatable}{lemma}{blateqci}~\blTitle{bl.ateqci}~ Equality of @ for values, inference.
\begin{align}
\vdash \tau_1 \in \mathbf{T}\quad\&\quad \vdash \tau_2 \in
    \mathbf{T}\quad\&\quad \vdash \mc{A} = \mc{B}\quad\&\quad \vdash \tau_1 =
    \tau_2\quad\Rightarrow\quad \vdash \mathbf{I} \lsem ( \mc{A} @ \tau_1 )
    \rsem = \mathbf{I} \lsem ( \mc{B} @ \tau_2 )
    \rsem\label{eq:bl.ateqci}\tag{bl.ateqci}
\end{align}
\end{restatable}

\begin{proof}
\begin{align}
  1 &&  & \vdash \tau_1 = \tau_2&\text{Hyp~ateqci.4}\notag%bl.ateqci.4
  \\ 2 &&  & \vdash \mc{A} = \mc{B}&\text{Hyp~ateqci.3}\notag%bl.ateqci.3
  \\ 3 &&  & \vdash \tau_1 \in \mathbf{T}&\text{Hyp~ateqci.1}\notag%bl.ateqci.1
  \\ 4 &&  & \vdash \tau_2 \in \mathbf{T}&\text{Hyp~ateqci.2}\notag%bl.ateqci.2
  \\ 5 &&  & \vdash ( ( \tau_1 = \tau_2 \wedge \mc{A} = \mc{B} ) \rightarrow
      \mathbf{I} \lsem ( \mc{A} @ \tau_1 ) \rsem =
      \mathbf{I} \lsem ( \mc{B} @ \tau_2 ) \rsem
      )&(\mbox{\ref{eq:df-bl.ateqc}},3,4)\notag
  \\ 6 &&  & \vdash \mathbf{I} \lsem ( \mc{A} @ \tau_1 ) \rsem = \mathbf{I}
      \lsem ( \mc{B} @ \tau_2 )
      \rsem&(\mbox{\ref{eq:mp2an}},1,2,5)\notag
\end{align}
\end{proof}

\begin{restatable}{lemma}{blatintroc}~\blTitle{bl.atintroc}~ @ Introduction: if a value is constant
its value is the same at any time.
\begin{align}
\vdash \tau \in \mathbf{T}\quad\Rightarrow\quad \vdash ( \mathbf{I} \lsem
    \mc{A} \rsem = \mc{B} \rightarrow \mathbf{I} \lsem ( \mc{A} @ \tau ) \rsem
    = \mc{B} )\label{eq:bl.atintroc}\tag{bl.atintroc}
\end{align}
\end{restatable}

\begin{proof}
\begin{align}
  1 &&  & \vdash \tau \in \mathbf{T}&\text{Hyp~atintroc.1}\notag%bl.atintroc.1
  \\ 2 &&  & \vdash \mathbf{I} \lsem ( \mc{A} @ \tau ) \rsem = \mathbf{I} \lsem
      \mc{A} , \tau
      \rsem&(\mbox{\ref{eq:df-bl.atc}},1)\notag
  \\ 3 &&  & \vdash \tau \in
      \mathbf{T}&\text{Hyp~atintroc.1}\notag%bl.atintroc.1
  \\ 4 &&  & \vdash ( \mathbf{I} \lsem \mc{A} \rsem = \mc{B} \rightarrow
      \mathbf{I} \lsem \mc{A} , \tau \rsem = \mc{B}
      )&(\mbox{\ref{eq:ax-bl.modelsc}},3)\notag
  \\ 5 &&  & \vdash ( \mathbf{I} \lsem \mc{A} \rsem = \mc{B} \rightarrow
      \mathbf{I} \lsem ( \mc{A} @ \tau ) \rsem = \mc{B}
      )&(\mbox{\ref{eq:syl5eq}},2,4)\notag
\end{align}
\end{proof}

\begin{restatable}{metadefinition}{dfblatrt}~\blTitle{df-bl.atrt}~ Applying @ retains type.
\begin{align}
\vdash \tau \in \mathbf{T}\quad\Rightarrow\quad \vdash ( \mc{A} \in \mc{B}
    \rightarrow \mathbf{I} \lsem ( \mc{A} @ \tau ) \rsem \in \mc{B}
    )\label{eq:df-bl.atrt}\tag{df-bl.atrt}
\end{align}
\end{restatable}

